\section{POSTPONED PROOFS}
\label{sec:postponed_proofs}
Here we present the proof of all results we skipped in the Section \ref{sec:main_results}.
% \subsection{Game Mechanics and Flowchart}
% \label{app:flowchart}

% The sequential interaction between Nature, the Sender, and the Receiver is summarized in Fig. \ref{fig:game_mechanics}. This vertical flow highlights the two-stage decision process and the noisy channel transmission.

% \begin{figure}[htbp]
%   \centering
%   \begin{tikzpicture}[
%       node distance=1.1cm,
%       block/.style={rectangle, draw, fill=blue!5, text width=2.1cm, text centered, rounded corners, minimum height=0.6cm, font=\scriptsize},
%       decision/.style={diamond, draw, fill=orange!5, text width=1.4cm, text centered, inner sep=0pt, minimum height=0.6cm, font=\scriptsize},
%       loss/.style={ellipse, draw, fill=red!5, text width=1.6cm, text centered, minimum height=0.5cm, font=\tiny},
%       line/.style={draw, ->, thick},
%       >=stealth
%   ]
%     % Nature
%     \node [block] (nature) {\textbf{Nature}\\ $(u,x,y)$};
    
%     % Sender
%     \node [block, below=of nature] (sender) {\textbf{Sender}\\ receives $(u,x)$};
%     \node [decision, below=of sender] (s_dec) {$A$};
    
%     % Noisy Channel
%     \node [block, below=of s_dec] (channel) {\textbf{Channel}\\ message $M$};
    
%     % Receiver Decision
%     \node [block, below=of channel] (receiver) {\textbf{Receiver}\\ receives $M$};
%     \node [decision, below=of receiver] (r_dec) {$\hat{y}$};

%     % Outcomes / Feedback
%     \node [loss, below left=0.4cm and 0.1cm of r_dec.center] (sl) {Sender Loss\\ $\ell(u, \hat{y})$};
%     \node [loss, below right=0.4cm and 0.1cm of r_dec.center] (rl) {Receiver Loss\\ $\ell(y, \hat{y})$};

%     % Connections
%     \path [line] (nature) -- node[right, font=\tiny] {$(u,x)$} (sender);
%     \path [line] (sender) -- (s_dec);
    
%     % Sender choices
%     \draw [line] (s_dec.west) -- ++(-0.4,0) |- node[left, pos=0.25, font=\tiny] {$\emptyset$} (channel.west);
%     \draw [line] (s_dec.east) -- ++(0.4,0) |- node[right, pos=0.25, font=\tiny] {$x$} (channel.east);
    
%     \path [line] (channel) -- node[right, font=\tiny] {$M$} (receiver);
%     \path [line] (receiver) -- (r_dec);
    
%     % Receiver choices
%     \draw [line] (r_dec.west) -- (sl.north);
%     \draw [line] (r_dec.east) -- (rl.north);
    
%     % Evaluation paths
%     \path [line, dashed, thin] (nature.west) -- ++(-1.1,0) |- node[left, near start, font=\tiny] {$u$} (sl.west);
%     \path [line, dashed, thin] (nature.east) -- ++(1.1,0) |- node[right, near start, font=\tiny] {$y$} (rl.east);

%   \end{tikzpicture}
%   \vspace{0.1cm}
%   \caption{Sequential game flowchart illustrating the timing of moves and loss feedback.}
%   \label{fig:game_mechanics}
% \end{figure}

\subsection{Proofs for Strategic Advantage Condition for Regression}
\label{app:regression}

This section provides the complete derivations for the general regression risk formulations and the subsequent sufficiency and necessity conditions.

\begin{proof}[Proof of Lemma~\ref{lem:rv_regression}]
We compute the risk in the mandatory disclosure setting for a linear predictor $b \in \R^d$ with $\yhatvemt=0$ and $\yhatvx = b^\top x$. The risk is given by:
\begin{align*}
  R_V(b, q) &= q \E[Y^2] + (1-q) \E[(Y - b^\top X)^2].
\end{align*}

Since all variables have mean $0$, $\Var(Y) = \E[Y^2]$. We now examine the term $\E[(Y-b^\top X)^2]$. We add and subtract $b_Y^\top X$ inside the expectation, where $b_Y = \Var(X)^{-1}\Cov(X,Y)$:
\begin{align*}
  \E[(Y - b^\top X)^2] &= \E[(Y - b_Y^\top X)^2] \\
  &+ \E[((b_Y - b)^\top X)^2] \\
  &+ 2\E[(Y - b_Y^\top X)(b_Y - b)^\top X].
\end{align*}

By the orthogonality principle of OLS, $\E[(Y - b_Y^\top X)X] = 0$, so the cross-term vanishes. Thus,
\begin{align*}
  \E[(Y - b^\top X)^2] &= \Var(Y) - \Var(b_Y^\top X) + \Var((b-b_Y)^\top X).
\end{align*}

Substituting back into the risk yields \eqref{eq:rv_b}. The optimal strategy minimizes $\Var((b-b_Y)^\top X)$ by choosing $b = b_Y$.
\end{proof}

\begin{proof}[Proof of Lemma~\ref{lem:sender_response}]
Given a realization $(x, u)$, the Sender selects the action $A \in \{\emptyset, x\}$ minimizing $(u - \yhatm)^2$. The expected loss from revealing is $q\,u^2 + (1-q)(u - b^\top x)^2$, while from withholding it is $u^2$. The Sender prefers revealing if:
\begin{align*}
  u^2 &> q\,u^2 + (1-q)(u - b^\top x)^2.
\end{align*}
Simplifying yields $(b^\top x)(2u - b^\top x) > 0$, which is $F(x,u) > 0$.
\end{proof}

\begin{proof}[Proof of Lemma~\ref{lem:rs_regression}]
For a fixed parameter $b$, the strategy of the Receiver is $\yhatvemt=0$ for an empty message $M=\emptyset$ and $\yhatvx=b^\top x$ for a successfully received feature vector $M=X$. 

According to Lemma~\ref{lem:sender_response}, the Sender reveals the feature vector $x$ if and only if $F(X,U) > 0$. 
Using the law of total expectation, we decompose the strategic risk $R_S(b, q) = \E[\ell(Y, \hat{y}(M))]$ into the two cases of $F$:
\begin{align*}
  R_S(b, q) &= \E\big[ \Ind{F \le 0} \ell(Y, \hat{y}(\emptyset)) \\
  & + \Ind{F > 0} ( q \ell(Y, \hat{y}(\emptyset)) + (1-q) \ell(Y, \hat{y}(X)) ) \big] \\
  &= \E\big[ \Ind{F \le 0} Y^2\\
  &+ \Ind{F > 0} ( q Y^2 + (1-q) (Y - b^\top X)^2 ) \big].
\end{align*}

Applying the identity $\Ind{F \le 0} = 1 - \Ind{F > 0}$:
\begin{align*}
  R_S(b, q) &= \E [ Y^2 ] - (1-q) \E \big[ \Ind{F > 0} ( Y^2 - (Y - b^\top X)^2 ) \big] \\
  &= \Var(Y) - (1-q) \E \big[ \Ind{F > 0} ( 2Y (b^\top X) \\
  &- (b^\top X)^2 ) \big],
\end{align*}
where $\Var(Y) = \E[Y^2]$ follows from the zero-mean assumption. 
Substituting $G(X,Y) = (2Y - b^\top X) b^\top X$ yields:
\begin{equation*}
  R_S(b, q) = \Var(Y) - (1-q) \E [ \Ind{F > 0} G(X,Y) ],
\end{equation*}
which concludes the proof.
\end{proof}

\begin{proof}[Proof of Theorem~\ref{thm:sufficiency_reg} (Sufficiency Condition)]
We seek condition such that $R_S^\star(q) - R_V^\star(q) < 0$. From Lemmas~\ref{lem:rv_regression} and~\ref{lem:rs_regression}:
\begin{align*}
  R_S^\star(q) - R_V^\star(q) &= \Var(Y) - (1-q)\E[G\Ind{F > 0}] \\
  &\quad - \big(\Var(Y) - (1-q)\Var(b_Y^\top X)\big).
\end{align*}
Using $\Ind{F > 0} = 1 - \Ind{F \le 0}$, we have $\E[G\Ind{F > 0}] = \E[G] - \E[G\Ind{F \le 0}]$.

To compute $\E[G]$, recall $G(X,Y) = (2Y - b^\top X)b^\top X$. Since $\E[X]=0$, $\Var(b^\top X) = \E[(b^\top X)^2]$. Furthermore, the optimal OLS weights $b_Y = \E[XX^\top]^{-1}\E[XY]$ satisfy $\E[Y(b^\top X)] = b^\top \E[XY] = b^\top \Var(X)b_Y = \Cov(b^\top X, b_Y^\top X)$. Thus,
\begin{align*}
  \E[G] &= 2\E[Y(b^\top X)] - \E[(b^\top X)^2] \\
  &= 2b^\top \Var(X)b_Y - \Var(b^\top X).
\end{align*}
From the identity $\Var((b-b_Y)^\top X) = \Var(b^\top X) + \Var(b_Y^\top X) - 2b^\top \Var(X)b_Y$, we obtain:
\begin{equation*}
  \E[G] = \Var(b_Y^\top X) - \Var((b{-}b_Y)^\top X).
\end{equation*}

Substituting this into the risk difference we get:
\begin{equation*}
  R_S^\star(q) - R_V^\star(q) = (1{-}q)\big( \Var((b{-}b_Y)^\top X) + \E[G\Ind{F \le 0}] \big).
\end{equation*}

We upper bound $\E[G\Ind{F \le 0}]$ by decomposing:
\begin{align*}
  \E[G\Ind{F \le 0}] &= \E[G\Ind{G \le 0}] \\
  &\quad + \E[G(\Ind{F \le 0} - \Ind{G \le 0})] \\
  &= -\tfrac{1}{2}\big(\E[|G|] - \E[G]\big) \\
  &\quad + \E[G \Ind{\sgn(F) \neq \sgn(G)}].
\end{align*}
When $\sgn(F) \neq \sgn(G)$, we have $|G| \le |G - F|= 2|Y{-}U||b^\top X|$. By Cauchy-Schwarz:
\begin{align*}
  \E[G \Ind{\sgn(F) \neq \sgn(G)}] &\le \E[2|Y{-}U|\,|b^\top X|] \\
  &\le 2\|U{-}Y\|_2 \sqrt{\Var(b_Y^\top X)}.
\end{align*}
Setting upper bound to be negative, (which would imply $R_S^\star(q) - R_V^\star(q) <0$), we get:
\begin{align*}
  &\Var((b{-}b_Y)^\top X) \\
  &- \tfrac{1}{2}\big( \E[|G|] - \Var(b_Y^\top X) \\
  &\quad + \Var((b{-}b_Y)^\top X) \big) \\
  &+ 2\|U{-}Y\|_2\sqrt{\Var(b_Y^\top X)} < 0.
\end{align*}
Rearranging isolates $\|U{-}Y\|_2$, yielding \eqref{eq:sufficiency_reg}.
\end{proof}

\begin{proof}[Proof of Theorem~\ref{thm:necessary_reg} (Necessary Condition)]
We prove the contrapositive: if \eqref{eq:necessary_reg} fails, then no strategic advantage can exist, i.e., $R_S(b,q) \ge R_V^\star(q)$ for all $b$.
From Lemma~\ref{lem:rv_regression} and Lemma~\ref{lem:rs_regression}, we have:
\begin{equation*}
  R_S(b,q) - R_V^\star(q) = (1{-}q)\big( \Var((b{-}b_Y)^\top X) + \E[G\Ind{F \le 0}] \big).
\end{equation*}

A strategic advantage requires $R_S(b,q) - R_V^\star(q) < 0$ for some $b$. We can lower bound the term $\E[G\Ind{F \le 0}]$ purely algebraically. 

If $G \le 0$, we have $G\Ind{F \le 0} \ge G = G\Ind{G \le 0}$. If $G > 0$, we have $G\Ind{F \le 0} \ge 0 = G\Ind{G \le 0}$. Therefore, pointwise $G\Ind{F \le 0} \ge G\Ind{G \le 0}$, which implies:
\begin{equation*}
  \E[G\Ind{F \le 0}] \ge \E[G\Ind{G \le 0}].
\end{equation*}

Substituting this inequality yields:
\begin{equation*}
  R_S(b,q) - R_V^\star(q) \ge (1{-}q)\big(\Var((b{-}b_Y)^\top X) + \E[G\Ind{G \le 0}]\big).
\end{equation*}

Now $\E[G\Ind{G \le 0}]$ depends only on $(X, Y)$ and $b$, with $U$ entirely eliminated. Using the identity $\E[G\Ind{G \le 0}] = -\frac{1}{2}(\E[|G|] - \E[G])$ and $\E[G] = \Var(b_Y^\top X) - \Var((b{-}b_Y)^\top X)$:
\begin{align*}
  R_S(b,q) - R_V^\star(q) &\ge (1{-}q)\Big(\Var((b{-}b_Y)^\top X) \\
  &\quad - \tfrac{1}{2}\big(\E[|G|] - \Var(b_Y^\top X) \\
  &\quad + \Var((b{-}b_Y)^\top X)\big)\Big).
\end{align*}

Multiplying by 2 and rearranging, $R_S(b,q) - R_V^\star(q) < 0$ requires:
\begin{equation*}
  \E[|G|] - \Var(b_Y^\top X) - \Var((b{-}b_Y)^\top X) > 0.
\end{equation*}

If this fails for all $b$, then $R_S(b,q) - R_V^\star(q) \ge 0$ for all $b$, meaning no strategic advantage is possible. Hence the necessary condition is that there exists some $b$ satisfying \eqref{eq:necessary_reg} and thus 
\begin{equation*}
  \max_b\E[|G|] - \Var(b_Y^\top X) - \Var((b{-}b_Y)^\top X) > 0.
\end{equation*}
\end{proof}

\subsection{Proofs for Strategic Advantage Condition for Classification}
\label{app:classification}

\begin{proof}[Proof of Lemma~\ref{lem:rv_classification}]
Under the mandatory disclosure setup (vanilla benchmark), the Receiver implements $\yhatvemt = \sgn(\E[Y])$ when $M = \emptyset$, and $\yhatvx = \sgn(\E[Y \mid X=x])$ when $M = X$.
Let $\eta(x) = \E[Y \mid X=x]$ and $p(x) = \Prob(Y=1 \mid X=x)$. Since $Y \in \{-1, +1\}$, we have $\eta(x) = 2p(x) - 1$.
When $\eta(x) \ge 0$, the optimum prediction is $+1$, yielding an error probability of $1 - p(x)$. When $\eta(x) < 0$, the optimum prediction is $-1$, yielding an error probability of $p(x)$. In both cases, the conditional Bayes risk is equivalently given by $(1 - |\eta(x)|)/2$. Thus, the overall expected risk when $M=X$ is $\E[(1 - |\eta(X)|)/2]$.

When $M = \emptyset$, the constant classifier yields risk $\Prob(Y \neq \sgn(\E[Y]))$, which is equal to $\Prob(Y = -\sgn(\E[Y]))$. 
Combining these mutually exclusive events with their respective probabilities $q$ and $1-q$ under the noisy channel yields the stated vanilla risk in \eqref{eq:rv_classification}.
\end{proof}

\begin{proof}[Proof of Lemma~\ref{lem:rs_classification}]
For a chosen default prediction $\yhatm = i$ when $M = \emptyset$, the Receiver correctly learns from a non-empty message $M = X$ that $U = -i$. Consequently, the optimal conditional strategy becomes $\yhatvx = \sgn(\E[Y \mid X=x, U=-i])$. 
If the Sender has preference $U = i$, withholding features gives a loss of $0$ (since the Receiver plays $i$), while revealing may risk a different prediction. Thus, the Sender strictly withholds the document, causing $M = \emptyset$. The conditional risk for this group ($U = i$) is simply the risk of the default prediction: $\Prob(Y = -i \mid U=i)$.
If the Sender has preference $U = -i$, withholding yields the worst possible loss for them. Hence, they choose to reveal $A = X$. The message drops with probability $q$ (leading to prediction $i$ and contributing $q\Prob(Y = -i \mid U=-i)$ to the error), and successfully arrives with probability $1-q$, where the conditional Bayes risk over the subgroup is $\E[(1-|\eta_{-i}(X)|)/2]$.
Multiplying the expectations by their respective cluster probabilities $\pi_i = \Prob(U=i)$ exactly formulates the overall strategic risk of $\cS_i$ as \eqref{eq:rs_classification}.
\end{proof}

\begin{proof}[Proof of Theorem~\ref{thm:advantage_classification}]
The risk functions $R_{\cS_1}^\star(q)$ and $R_{\cS_{-1}}^\star(q)$ from Lemma~\ref{lem:rs_classification} are linear equations in $q$. Assuming $\E[Y] > 0$, which implies $\Prob(Y=1) > \Prob(Y=-1)$. At maximum channel noise ($q=1$), the messages are completely dropped ($M = \emptyset$), reducing the Receiver risk to the static default prediction's baseline error. Hence, $R_{\cS_{1}}^\star(1) = \Prob(Y=-1)$, and $R_{\cS_{-1}}^\star(1) = \Prob(Y=1)$. Since $\Prob(Y=1) > \Prob(Y=-1)$, it is clear that $R_{\cS_{1}}^\star(1) < R_{\cS_{-1}}^\star(1)$, making $\cS_1$ strictly optimal at $q=1$. Because both risks are lines as functions of $q$, if $R_{\cS_{1}}^\star(0) \le R_{\cS_{-1}}^\star(0)$ (lets call it scenarion A), $\cS_{1}$ is optimal at both extremes and thus dominates everywhere via linear interpolation. Conversely, if $R_{\cS_{-1}}^\star(0) < R_{\cS_{1}}^\star(0)$ (lets call it scenario B), the linear functions cross once inside the interval $q \in (0,1)$. This intersection naturally defines the transition threshold $q^\star$ given by 
    \begin{align}
    q^\star = \frac{R_{\cS_1}^\star(0) - R_{\cS_{-1}}^\star(0)}{(R_{\cS_{-1}}^\star(1) - R_{\cS_{-1}}^\star(0)) - (R_{\cS_1}^\star(1) - R_{\cS_1}^\star(0))},
    \end{align}
such that for $q < q^\star$, $\cS_{-1}$ is optimal, and for $q \ge q^\star$, $\cS_{1}$ is optimal. We obtain the strategic advantage condition by verifying when $R_V^\star(q) - R_{\cS_i}^\star(q) > 0$ for the optimal $i$. Under the domain where $\cS_1$ is Stackelberg-optimal (all $q$ in scenario A, or $q \ge q^\star$ in scenario B):
\begin{align*}
  R_V^\star(q) - R_{\cS_1}^\star(q) > 0.
\end{align*}
Using $\sgn(\E[Y]) = 1$, we decompose the baseline error rate in $R_V^\star(q)$:
\begin{align*}
  q\Prob(Y=-1) &= q [ \pi_1 \Prob(Y=-1 \mid U=1) \\
  &\quad + \pi_{-1} \Prob(Y=-1 \mid U=-1) ].
\end{align*}
Upon substituting $R_V^\star(q)$ and $R_{\cS_1}^\star(q)$, the terms matching $q\pi_{-1}\Prob(Y=-1 \mid U=-1)$ cancel algebraically. The remaining terms structurally share a common $(1-q)$ factor. Dividing the strict inequality by $(1-q)$ yields the noise-invariant condition shown in \eqref{eq:adv_case1}. In the alternate domain where $\cS_{-1}$ is strictly optimal ($q < q^\star$ in scenario B), we calculate $R_V^\star(q) - R_{\cS_{-1}}^\star(q) > 0$. Here, the $q$ dependencies no longer cancel out purely because the vanilla baseline ($+1$) and the strategic default choice ($-1$) are misaligned. Grouping the independent expectations corresponding to $q$ and $(1-q)$ distinctly and isolating terms directly formulates the threshold specified in \eqref{eq:adv_case2}.
\end{proof}

\subsection{Online Learning Analysis}
\label{app:online_learning}

\begin{proof}[Proof of Theorem~\ref{thm:online_regret} (Outline)]
We decompose the cumulative regret into two components and bound each separately.

\textbf{Step 1: Regret Decomposition.}
Let $D^\star \in \{\text{Mandatory}, \text{Voluntary}\}$ denote the optimal disclosure regime and $\hat{y}^\star$ the associated optimal predictor under the true environment $\theta^\star$. The instantaneous regret at round $t$ is:
\begin{align*}
  r_t &= \ell(y_t, \hat{y}_t(M_t)) - \ell(y_t, \hat{y}^\star(M_t^\star)).
\end{align*}
We decompose this as:
\begin{align*}
  r_t &= \underbrace{\ell(y_t, \hat{y}_t(M_t)) - \ell(y_t, \hat{y}^\star_t(M_t))}_{\text{(I): estimation regret}} \\
  &\quad + \underbrace{\ell(y_t, \hat{y}^\star_t(M_t)) - \ell(y_t, \hat{y}^\star(M_t^\star))}_{\text{(II): regime selection regret}},
\end{align*}
where $\hat{y}^\star_t$ denotes the oracle-optimal predictor for the regime $D_t$ chosen at round $t$.

\textbf{Step 2: Bounding Regime Selection Regret (II).}
Term (II) is nonzero only when the algorithm selects a suboptimal regime $D_t \neq D^\star$. Define the suboptimality gap $\Delta = |R_{D^\star}^\star - R_{D^c}^\star|$ where $D^c$ is the suboptimal regime. Under the Bayesian posterior update, posterior concentration yields:
\begin{equation*}
  \Prob(D_t \neq D^\star) \le \Prob\big(\|\hat{\theta}_t - \theta^\star\| > \delta\big),
\end{equation*}
for an appropriate $\delta$ depending on $\Delta$. By standard Bayesian posterior contraction under the compactness and regularity assumptions on $\Theta$, the posterior concentrates at rate $1/\sqrt{t}$:
\begin{equation*}
  \Prob\big(\|\hat{\theta}_t - \theta^\star\| > \delta\big) \le C_1 \exp\big(-c_1 t \delta^2\big).
\end{equation*}
Summing over $t$, the total contribution of regime misidentification rounds is:
\begin{equation*}
  \sum_{t=1}^T \E[\text{(II)}_t] \le C_2 \log T.
\end{equation*}

\textbf{Step 3: Bounding Estimation Regret (I).}
Conditional on selecting the correct regime $D_t = D^\star$, term (I) reduces to standard online prediction regret. Under the mandatory regime, this is online linear regression against an i.i.d.\ sequence, which satisfies $\sum_t \E[\text{(I)}_t \mid D_t = D^\star] \le C_3 \sqrt{T}$ via standard results (e.g., Vovk--Azoury--Warmuth bounds for squared loss).

Under the voluntary regime, the Receiver additionally needs to estimate the Sender's best-response mapping. Since the Sender's action depends on the committed predictor and the joint distribution, the effective observation model remains parametric. The Bayesian belief update in Algorithm~\ref{alg:adaptive_disclosure} incorporates the likelihood $\Prob(M_t, y_t \mid \theta, \hat{y}_t)$, which accounts for the Sender's strategic behavior. The posterior mean predictor then achieves:
\begin{equation*}
  \sum_{t=1}^T \E[\text{(I)}_t] \le C_3 \sqrt{T}.
\end{equation*}

\textbf{Step 4: Combining.}
Adding both terms:
\begin{equation*}
  \E[\mathcal{R}_T] = \sum_{t=1}^T \E[r_t] \le C_2 \log T + C_3 \sqrt{T} \le C\sqrt{T},
\end{equation*}
where the $\log T$ regime identification cost is absorbed into the dominant $\sqrt{T}$ estimation term.
\end{proof}
